\documentclass[11pt]{article}

\usepackage[margin=25mm]{geometry}

\usepackage{amsmath,amssymb,fontspec,unicode-math}

\IfFontExistsTF{DejaVu Serif}{\setmainfont{DejaVu Serif}}{\setmainfont{Latin Modern Roman}}

\IfFontExistsTF{Latin Modern Math}{\setmathfont{Latin Modern Math}}{\setmathfont{DejaVu Math TeX Gyre}}

\usepackage{xurl,hyperref}

\setlength{\parindent}{0pt}\setlength{\parskip}{6pt}

\begin{document}

\title{Provably Efficient Multi-Task Meta Bandit Learning via Shared Representations}

\author{ELVIA preliminary AI analysis}\date{21 October 2026}\maketitle

Preliminary assessment: Major.\par

Independent scientific review remains pending. Author adoption has not been requested.

The central source-regret and representation-estimation proofs use unjustified probabilistic steps. Target confidence, potential and excitation arguments also have explicit local failures. Separate released-notebook discrepancies do not establish that published figure rankings are false.

\section*{Review scope}

Claim-focused reading of the original paper and relevant author-released passages; direct mathematical or static-source reasoning. No author program, earlier verifier, model or benchmark was executed during this check.

\section*{Problem 1: Optimizing Gaussian candidates does not preserve centering}

Source: Theorem 4.3 and Remark 4.4, PDF pp.5–6; Appendix D, PDF pp.19–21

Assessment: supported; impact: major

\subsection*{Source claim (paraphrase)}

The regret proof treats oracle rewards as centered Gaussian martingale increments, and the task-scaling remark suppresses the largest singular value.

\subsection*{Analysis}

For independent candidates \(Z_1,Z_2\sim N(0,1)\) and parameter \(1\), the exploratory first candidate is centered but the oracle is \(\max(Z_1,Z_2)=(Z_1+Z_2+|Z_1-Z_2|)/2\), with mean \(1/\sqrt{\pi}\). Thus exploration regret has positive drift \(N_1/\sqrt{\pi}\), and the asserted martingale implication fails. On \(Z_1<0\), even the preceding pointwise inequality \(\max(Z_1,Z_2)-Z_1\le\max(Z_1,Z_2)\) fails. Separately, rank-one task weights all equal to \(1\) satisfy bounded task norms and incoherence but give \(\sigma_{\max}^{\star}=\sqrt T\); the retained spectral factor cannot be absorbed into a task-independent constant.

\subsection*{Finding and scope}

The central proof steps and simplified task dependence are unsupported. This contextual candidate example isolates selection bias; it is not claimed to instantiate every source-budget premise or to disprove every possible corrected regret theorem.

\subsection*{Proposed correction}

Under an added bounded-action model, valid starting bounds are \(R_{\rm explore}\le2LuTN_1\) and \(R_{\rm commit}\le2L(N-N_1)\sum_t\|\widehat\theta_t-\theta_t^\star\|\). A coherent candidate-generation model and a new regret analysis are still required.

A proposed correction is retained where stated; independent verification and author adoption remain unestablished.

\section*{Problem 2: A fitted representation cannot be pulled outside the same-sample expectation}

Source: Algorithm 1, PDF p.4; Lemma C.1, PDF p.17

Assessment: supported; impact: major

\subsection*{Source claim (paraphrase)}

The proof substitutes \(\mathbb E[\widehat B^\top xx^\top\widehat B]=I\) and independent Gaussian reduced-design summands although the same samples fit \(\widehat B\).

\subsection*{Analysis}

Let \(d=8,r=N_1=T=1,\theta^\star=e_1\), zero noise and \(x\sim N(0,I)\). The algorithm retains the lone observation \(y=x_1\) and obtains \(\widehat B=\operatorname{sign}(x_1)x/\|x\|\). Consequently \((\widehat B^\top x)^2=\|x\|^2\) has expectation \(8\), rather than \(1\). An independent fresh \(x^\prime\), conditional on the fitted basis, does have expectation \(1\).

\subsection*{Finding and scope}

The displayed expectation/independence justification is invalid. The small diagnostic does not meet Theorem 4.1’s high-sample conditions and is not a complete theorem counterexample; a dependence-aware proof remains necessary.

\subsection*{Proposed correction}

Independent representation-fitting and regression blocks justify the conditional Gaussian calculation, but change sample accounting. Preserving sample reuse requires a uniform or dependence-aware concentration argument.

A proposed correction is retained where stated; independent verification and author adoption remain unestablished.

\section*{Problem 3: The rank-one confidence simplification drops a logarithmic coefficient}

Source: Theorem 5.1, PDF p.7; Appendix E, PDF p.24

Assessment: supported; impact: major

\subsection*{Source claim (paraphrase)}

The determinant radius squared contains \(\sigma^2[\log(\det V_n/\lambda^r)+2\log(1/\delta)]\), while the simplification has \(r\log(1/\delta)\).

\subsection*{Analysis}

At \(r=1\), grant a known representation \(e_1\) in \(d=100\), \(x_m=e_1,\theta^\star=0,\lambda=L=\sigma=1,n=1000,S=0,\delta=10^{-12}\). Independent unit Gaussian noise gives design-norm error \(\sqrt{1000/1001}|Z|\). The printed radius \(\sqrt{\log(1001/10^{-12})}=5.877055\) has exact two-sided failure probability \(4.10284256\times10^{-9}\), exceeding \(\delta+2d^{-10}\approx10^{-12}\). This tests the conditional target calculation, bypassing rather than authenticating source-learning premises.

\subsection*{Finding and scope}

The rank-one simplification is not a valid confidence bound on its stated conditional domain. For ranks at least two, the displayed coefficient comparison alone is not an objection; the determinant form remains the proper starting point.

\subsection*{Proposed correction}

Preserve \(2\log(1/\delta)\); the trace bound gives noise radius \[\sigma\sqrt{r\log(1+nL^2/(r\lambda))+2\log(1/\delta)}\]. Retain the stated representation-bias terms and recheck dependent regret results. The compact-expression issue may be inherited from the cited OFUL result; no priority claim is made.

A proposed correction is retained where stated; independent verification and author adoption remain unestablished.

\section*{Problem 4: The untruncated potential inequality needs a size condition}

Source: Appendix F, Eqs.(12)–(13), PDF p.25; Section 5 assumptions

Assessment: supported; impact: major

\subsection*{Source claim (paraphrase)}

The proof uses \(\sum_m x_m^\top\overline V_{m-1}^{-1}x_m\le2\log(\det\overline V_N/\lambda^d)\) for unrestricted \(\lambda>0,L>0\).

\subsection*{Analysis}

With \(d=3,N=1,\lambda=1,x_1=(1000,0,0)^\top,L=1000\), the left side is \(10^6\), but the right side is \(2\log(1000001)=27.6310231\). The generally valid inequality truncates each quadratic form at \(1\). Dropping truncation is justified, for example, by \(\lambda\ge L^2\), which is absent here.

\subsection*{Finding and scope}

This deterministic step used in the target-regret proof is false on its stated domain; it does not by itself disprove every asymptotic rate under a normalized-action convention.

\subsection*{Proposed correction}

Impose the needed size condition, or keep \(a=L^2/\lambda\) and factor \(c(a)=a/\log(1+a)\). Since \(q/\log(1+q)\) increases, \(\sum q_m\le c(a)\log(\det\overline V_N/\lambda^d)\). Propagate the factor through a valid conditional confidence argument.

A proposed correction is retained where stated; independent verification and author adoption remain unestablished.

\section*{Problem 5: Target sample count does not ensure excitation}

Source: Theorem G.1 and Algorithm 3, PDF pp.25–26; source-only Assumption 3.1

Assessment: supported; impact: major

\subsection*{Source claim (paraphrase)}

The target regression theorem uses an ordinary reduced-design inverse and only a target sample-count lower bound; its Gaussian action assumption is explicitly for source tasks.

\subsection*{Analysis}

Grant \(\widehat B=B^\star=e_1\in\mathbb R^2,\theta^\star=e_1\), zero target noise and target observations \(x_m=e_2\). For every sample count, \(\Phi\widehat B=0,Y=0\). The displayed pseudoinverse estimate is zero, with error \(1\), versus claimed bound \(1.12(0.01)=0.0112\). Greedy collection from ordered actions \(\{e_2,e_1\}\) with zero initial estimate and a first-index tie rule can remain unexcited throughout.

\subsection*{Finding and scope}

The target-design argument omits identification/excitation. This known-representation construction isolates the target obstruction, without claiming a full source-process counterexample.

\subsection*{Proposed correction}

Add a justified target design condition such as \(\lambda_{\min}(\widehat B^\top\Phi^\top\Phi\widehat B)\ge\gamma N_2>0\), together with appropriate conditional noise and representation-residual assumptions. No complete repaired theorem is established.

A proposed correction is retained where stated; independent verification and author adoption remain unestablished.

\section*{Problem 6: The released generator uses variance as a standard deviation}

Source: Section 6, PDF pp.8–9; original Representation\_Learning/Synthetic/Algorithm and Transfer\_Learning/Synthetic/Algorithm notebooks

Assessment: supported; impact: minor

\subsection*{Source claim (paraphrase)}

The source notebook sets \(\sigma=10^{-3}\) and calls normal with scale \(\sigma^2\); the transfer notebook sets \(\sigma=10^{-1}\) and does the same.

\subsection*{Analysis}

Normal scale is standard deviation, so the actual variance is \(\sigma^4\): respectively \(10^{-12}\) and \(10^{-4}\), versus paper variances \(10^{-6}\) and \(10^{-2}\). This follows directly from \(\eta=\sigma^2Z\). Other notebooks using scale \(\sigma\) are countercontrols; the defect is not universal.

\subsection*{Finding and scope}

These settings instantiate different noise levels. Actual historical figure provenance and ranking changes remain unestablished.

\subsection*{Proposed correction}

When \(\sigma\) denotes standard deviation, pass scale \(\sigma\). Recompute affected comparisons before attributing performance at the stated variance.

A proposed correction is retained where stated; independent verification and author adoption remain unestablished.

\section*{Problem 7: The MoM comparison mixes decision clocks and action domains}

Source: Section 6.2, PDF p.9; Figure 1(a), p.10; synthetic source notebook, cell 7

Assessment: supported; impact: minor

\subsection*{Source claim (paraphrase)}

The comparison is described over \(200\) rounds, but the MoM branch has \(45\) initial, \(2\times45\) basis-direction and \(110\) commitment actions.

\subsection*{Analysis}

The literal total is \(245\), yielding \(246\) cumulative entries with the initial zero. Truncation to \(201\) entries retains only \(65\) commitment actions, while context indices \(45\)–\(89\) are reused in both basis loops. Probes \(\sqrt{\lambda_0}\widehat B_{:,k}\) also need not belong to the finite candidate set: \((e_1+e_2)/\sqrt2\notin\{e_1,e_2\}\). The notebook explicitly calls the baseline infinite-action, which supports an intentional variant but does not establish a common comparison protocol.

\subsection*{Finding and scope}

The released comparison lacks a single common clock/action domain. No claim is made that the original continuous-action MoM method is erroneous or a published ranking reverses.

\subsection*{Proposed correction}

Use one decision clock and explicitly align feasible sets and comparators. A changed benchmark requires new runs; array truncation alone does not repair it.

A proposed correction is retained where stated; independent verification and author adoption remain unestablished.

\section*{Problem 8: The MovieLens projector includes the full ambient basis}

Source: Transfer\_Learning/Movielens/Algorithm notebook, cells 9 and 11; subspace distance definition, PDF p.3

Assessment: supported; impact: minor

\subsection*{Source claim (paraphrase)}

For a \(5\times100\) rank-one task matrix, the notebook uses all economy-SVD left vectors in \(I-UU^\top\).

\subsection*{Analysis}

Repeat \((1,1,1,1,1)^\top/\sqrt5\) as all columns and take \(\widehat B=(1,-1,0,0,0)^\top/\sqrt2\). The true subspace distance is \(1\). Economy SVD nevertheless returns a square orthogonal \(U\), so \(UU^\top=I_5\) and the coded distance is zero for every estimate. The quantity enters the confidence radius; at \(n=100,L=S=\lambda=1\), erasing a unit error removes \(21\) from its bias term.

\subsection*{Finding and scope}

The local geometric measurement and radius change are established. A changed selected action or published MovieLens figure is not established; the notebook settings differ from the figure.

\subsection*{Proposed correction}

Use the intended signal-rank basis \(U_{:,1:r}\), then compute distance to that subspace. This fixes the projector, without certifying other confidence assumptions.

A proposed correction is retained where stated; independent verification and author adoption remain unestablished.

\section*{Limits}

Preliminary AI assessment; independent human scientific adjudication remains pending.\par

Analytic witnesses and static release checks do not remeasure the published benchmark results.\par

\section*{Source identity}

NeurIPS 2025 official proceedings PDF (retained local copy)

\url{https://proceedings.neurips.cc/paper_files/paper/2025/file/96d248deec8357a16f1fecda728d50c5-Paper-Conference.pdf}

Reviewed PDF SHA-256: \nolinkurl{362ae65b97a6c5de6ebfd6b6a65078b27d396ffb58653a49160fdeb4701de678}

The SHA-256 identifies the exact retained PDF inspected in this rapid source/argument check. Original supplied report bytes are retained locally. No unambiguous original PDF retrieval timestamp was recovered; file modification time was not substituted.

\end{document}
